% Generated draft: requires mathematical and source review.
\documentclass[11pt]{article}
\usepackage[margin=25mm]{geometry}
\usepackage{fontspec}
\setmainfont{Latin Modern Roman}
\usepackage{amsmath,amssymb,mathtools,unicode-math}
\setmathfont{Latin Modern Math}
\usepackage{enumitem,needspace,xurl,hyperref}
\setlength{\parindent}{0pt}
\setlength{\parskip}{5pt}
\setlength{\emergencystretch}{3em}
\newcommand{\N}{\mathbb{N}}
\newcommand{\Z}{\mathbb{Z}}
\newcommand{\Q}{\mathbb{Q}}
\newcommand{\R}{\mathbb{R}}
\newcommand{\C}{\mathbb{C}}
\newcommand{\F}{\mathbb{F}}
\newcommand{\A}{\mathbb{A}}
\newcommand{\PP}{\mathbb{P}}
\newcommand{\E}{\mathbb{E}}
\newcommand{\eps}{\varepsilon}
\newcommand{\dd}{\,\mathrm{d}}
\begin{document}
\section*{UnsolvedMath ID: 701719}
\textbf{Source problem code:} catalogue entry
\subsection{Definitions and mathematical statement}
For a finite simplicial graph \(\Gamma\), the associated right-angled Artin group is
\[
A_\Gamma=\left\langle v\in V(\Gamma)\ \middle|\ vw=wv\text{ whenever }\{v,w\}\in E(\Gamma)\right\rangle.
\]
Its Salvetti complex \(S_\Gamma\) has one vertex, one circle for each vertex of \(\Gamma\), and a cubical torus for each complete subgraph, attached along the corresponding lower-dimensional tori. Use the standard convention that a special cube complex is a nonpositively curved cube complex admitting a cubical local isometry to a Salvetti complex. A group \(H\) is \emph{virtually compact special} if there exist a finite-index subgroup \(H_0\leq H\) and a connected compact special cube complex \(X\) such that \(H_0\cong\pi_1(X)\).

For a group \(G\) and an element \(g\in G\), define its centralizer by
\[
C_G(g)=\{h\in G: hg=gh\}.
\]
The question proposed by A. Minasyan is whether, for every virtually compact special group \(G\) and every \(g\in G\), the group \(C_G(g)\) is virtually compact special. Explicitly, must there exist a subgroup \(K\leq C_G(g)\) and a connected compact special cube complex \(Y\) such that
\[
[C_G(g):K]<\infty
\qquad\text{and}\qquad
K\cong\pi_1(Y)?
\]
The quantifier over \(g\) includes the identity and elements of finite or infinite order; the supplied statement imposes no restriction on their order. The cube-complex terminology is expanded here using its standard meaning; the supplied record does not itself provide these definitions.
\subsection{Short English statement}
A. Minasyan asks whether the centralizer of every element of a virtually compact special group is again virtually compact special. Here this means having a finite-index subgroup isomorphic to the fundamental group of a connected compact special cube complex. The supplied registry labels the problem open; that label is not an independently verified assessment.
\subsection{Sources}
\begin{itemize}
\item Local UnsolvedMath registry, supplied complete source record: permanent ID \(701719\), exact source problem code \texttt{catalogue entry}, title \emph{Kourovka Notebook Problem 20.61}. The record gives A. Minasyan as proposer and \(2022\) as the proposed year. Its supplied status is \texttt{open}.
\item The supplied background attributes the problem to \emph{The Kourovka Notebook (Unsolved Problems in Group Theory)}, 20th issue (2022), and the 21st issue edition (2026), identified as \texttt{arXiv:1401.0300v48}; original problem number \(20.61\). The registry's collection description names E. I. Khukhro and V. D. Mazurov as editors.
\item Registry provenance: supplied upstream ID \(701719\), upstream revision \texttt{d371ee917ac9647fbb06cd6fc529deec6f3dda82}, and record update timestamp \texttt{2026-10-16T14:00:00Z}. These are supplied provenance data, not evidence of independent source or literature verification.
\end{itemize}
\end{document}
